% ----------------------------------------------------------------------
% Configurações visuais — Capacitação de MySQL
% ----------------------------------------------------------------------

% --- Paleta da marca ---------------------------------------------------
\definecolor{branco}{HTML}{FFFFFF}
\definecolor{roxoescuro}{HTML}{1E1647}
\definecolor{roxoclaro}{HTML}{614AD3}
\definecolor{amarelo}{HTML}{FFFF00}
\definecolor{cinzafundo}{HTML}{F4F4F7}
\definecolor{cinzalinha}{HTML}{9A9AB0}
\definecolor{verdeok}{HTML}{1B7F4B}
\definecolor{vermelhoerro}{HTML}{B3261E}

% --- Atalhos de marcação -----------------------------------------------
\newcommand{\comando}[1]{\texttt{#1}}
\newcommand{\sqlkw}[1]{\texttt{\textbf{#1}}}
\newcommand{\ok}{\textcolor{verdeok}{\ding{51}}}
\newcommand{\nok}{\textcolor{vermelhoerro}{\ding{55}}}

% ----------------------------------------------------------------------
% Ambientes de código (listings)
% ----------------------------------------------------------------------
% O 'literate' abaixo é necessário porque o listings não processa UTF-8
% sozinho: sem ele, cada acento dentro de um bloco de código vira lixo.
\lstset{
    inputencoding=utf8,
    extendedchars=true,
    literate=%
      {á}{{\'a}}1 {à}{{\`a}}1 {ã}{{\~a}}1 {â}{{\^a}}1
      {é}{{\'e}}1 {ê}{{\^e}}1 {í}{{\'i}}1
      {ó}{{\'o}}1 {õ}{{\~o}}1 {ô}{{\^o}}1
      {ú}{{\'u}}1 {ü}{{\"u}}1 {ç}{{\c c}}1
      {Á}{{\'A}}1 {À}{{\`A}}1 {Ã}{{\~A}}1 {Â}{{\^A}}1
      {É}{{\'E}}1 {Ê}{{\^E}}1 {Í}{{\'I}}1
      {Ó}{{\'O}}1 {Õ}{{\~O}}1 {Ô}{{\^O}}1
      {Ú}{{\'U}}1 {Ç}{{\c C}}1
      {º}{{\textordmasculine}}1 {ª}{{\textordfeminine}}1
}

\lstdefinestyle{base}{
    basicstyle       = \ttfamily\footnotesize,
    backgroundcolor  = \color{cinzafundo},
    frame            = single,
    rulecolor        = \color{cinzalinha},
    breaklines       = true,
    breakatwhitespace= false,
    postbreak        = \mbox{\textcolor{roxoclaro}{$\hookrightarrow$}\space},
    showstringspaces = false,
    upquote          = true,
    tabsize          = 2,
    captionpos       = b,
    xleftmargin      = 12pt,
    framexleftmargin = 12pt,
}

\lstdefinestyle{sql}{
    style        = base,
    language     = SQL,
    morekeywords = {AUTO_INCREMENT,ENGINE,InnoDB,UNSIGNED,IF,EXISTS,DATABASE,
                    REFERENCES,CONSTRAINT,CASCADE,RESTRICT,BOOLEAN,TINYINT,
                    DECIMAL,VARCHAR,TIMESTAMP,CHARACTER,COLLATE,LIMIT,SHOW,
                    DESCRIBE,TRUNCATE,USE,REPLACE},
    keywordstyle = \color{roxoescuro}\bfseries,
    commentstyle = \color{cinzalinha}\itshape,
    stringstyle  = \color{roxoclaro},
}

\lstdefinestyle{bash}{
    style        = base,
    language     = bash,
    keywordstyle = \color{roxoescuro}\bfseries,
    commentstyle = \color{cinzalinha}\itshape,
    stringstyle  = \color{roxoclaro},
}

\lstdefinestyle{yaml}{
    style         = base,
    basicstyle    = \ttfamily\footnotesize\color{roxoescuro},
    commentstyle  = \color{cinzalinha}\itshape,
    morecomment   = [l]{\#},
    moredelim     = [l][\color{black}]{:\ },
}

\lstnewenvironment{sqlcode}[1][]{\lstset{style=sql,#1}}{}
\lstnewenvironment{bashcode}[1][]{\lstset{style=bash,#1}}{}
\lstnewenvironment{yamlcode}[1][]{\lstset{style=yaml,#1}}{}

% ----------------------------------------------------------------------
% Caixas de destaque
% ----------------------------------------------------------------------
\newtcolorbox{aviso}[1][Atenção]{
    enhanced, breakable,
    colback   = vermelhoerro!5,
    colframe  = vermelhoerro,
    coltitle  = branco,
    fonttitle = \bfseries,
    title     = {#1},
    left=3mm, right=3mm, top=2mm, bottom=2mm,
}

\newtcolorbox{dica}[1][Dica]{
    enhanced, breakable,
    colback   = roxoclaro!5,
    colframe  = roxoclaro,
    coltitle  = branco,
    fonttitle = \bfseries,
    title     = {#1},
    left=3mm, right=3mm, top=2mm, bottom=2mm,
}

\newtcolorbox{conceito}[1][Conceito]{
    enhanced, breakable,
    colback   = roxoescuro!5,
    colframe  = roxoescuro,
    coltitle  = branco,
    fonttitle = \bfseries,
    title     = {#1},
    left=3mm, right=3mm, top=2mm, bottom=2mm,
}

% ----------------------------------------------------------------------
% Links e metadados do PDF
% ----------------------------------------------------------------------
\hypersetup{
    colorlinks  = true,
    linkcolor   = roxoescuro,
    citecolor   = roxoescuro,
    filecolor   = roxoclaro,
    urlcolor    = roxoclaro,
    pdftitle    = {Guia Prático para a Capacitação de MySQL},
    pdfauthor   = {Ronivaldo Domingues de Andrade},
    pdfsubject  = {Fundamentos de MySQL e Projetos de Bancos de Dados},
    pdfkeywords = {MySQL, SQL, Banco de Dados, SGBD, Modelo Relacional,
                   CRUD, JOIN, phpMyAdmin, Docker},
}
